\documentclass[10pt,a4paper]{amsart}
\usepackage[margin=19mm]{geometry}
\usepackage{amsmath,amssymb,mathtools}
\usepackage[scr=rsfs,cal=euler]{mathalfa}
\usepackage{xfrac}
\usepackage[unicode,hidelinks,bookmarks=false]{hyperref}
\newcommand{\ie}{i.e.\ }
\newcommand{\cf}{cf.\ }
\newcommand{\resp}{respectively\ }
\newcommand{\Subsection}{\subsection}
\providecommand{\nobreakdash}{\nobreak\hbox{-}}
\setlength{\emergencystretch}{2em}
\multlinegap0pt
\usepackage{bm}
\usepackage{bbm}
\usepackage{dsfont}
\usepackage{xspace}
\usepackage{cancel}
\usepackage{mathtools}
\usepackage{amsthm}
\makeatletter
\@ifundefined{theorem}{\newtheorem{theorem}{Theorem}}{}
\@ifundefined{lemma}{\newtheorem{lemma}[theorem]{Lemma}}{}
\makeatother
\title[Stability, bounded generation and strong boundedness]{Stability, bounded generation and strong boundedness}
\author{Alexander Trost}
\date{Source: arXiv:2305.11562v3 (CC BY 4.0)}
\begin{document}
\maketitle

\noindent\emph{Source and licence.} Stability, bounded generation and strong boundedness. Version arXiv:2305.11562v3 (CC BY 4.0), distributed under the Creative Commons Attribution 4.0 International licence, as recorded in the arXiv licence metadata for this exact version and shown on the abstract page https://arxiv.org/abs/2305.11562v3. Full rights notice: licenses/arxiv-2305.11562-CC-BY-4.0.txt.\par\smallskip

\noindent\textit{Editorial scope.} Counted unit: the Lemma whose source label is \texttt{fixing\_first\_column\_symplectic\_step2} in the article (source0.tex lines 624--626, source label \texttt{fixing\_first\_column\_symplectic\_step2}), delivered together with the article's own complete original proof of it (source0.tex lines 628--685, one \texttt{proof} environment of 58 lines). No mathematics is written, repaired or re-derived: statement and proof are the article's own text line for line. Required context retained uncounted: stability\_relative\_ideal (lines 551--553); stability\_sln\_prep\_step1 (lines 587--589); fixing\_first\_column\_symplectic\_step1 (lines 609--611). Every definition, display and label of those lines is retained unchanged. Omitted from this package and not redistributed: 1--550, 554--586, 590--608, 612--623, 627--627, 686--1144. Nothing inside a retained excerpt is omitted or altered: every retained body is delivered exactly as the article's lines are written, so no omission receipt of any kind accompanies this package. Citations: the retained excerpts carry no citation command at all, so no bibliography entry of the article is transcribed and no reference is redistributed. Reference adaptation: none was needed and none was made; every reference inside the delivered unit resolves inside it, so no unresolved reference or citation is produced. Printed slips kept literal: no printed word, symbol, hypothesis or number of the counted statement or of its proof was altered. Layout and class: the article's own class is \texttt{amsart} with packages including fullpage, babel, mathtools, amssymb, mathrsfs, nicefrac, upgreek, fontenc, graphicx, stmaryrd, booktabs, hyperref, xy, graphics; the wrapper is the lane's pinned \texttt{amsart} editorial wrapper at 10pt on A4, which restates the theorem-like environments the retained text needs and adds the article's own macro definitions that the retained text needs (none), with extra wrapper packages (bm, bbm, dsfont, xspace, cancel, mathtools). The delivered numbering is the unit's own. Assets: no retained span contains a figure environment, a graphics-inclusion command, a PDF-page inclusion, a table or a TikZ picture; no image file is copied into this package, the package declares no asset dependency and no image file is redistributed with it. Licence: version arXiv:2305.11562v3 of this article is distributed under the Creative Commons Attribution 4.0 International licence, as recorded in the arXiv licence metadata for this exact version and shown on the abstract page https://arxiv.org/abs/2305.11562v3. A full rights notice is delivered at licenses/arxiv-2305.11562-CC-BY-4.0.txt. Characters: every retained excerpt is pure ASCII, so the delivered engine, XeLaTeX with its default Latin Modern text font, reports no missing character.
\begin{lemma}\label{stability_relative_ideal}
Let $R$ be a Dedekind domain and $n\geq 2.$ Further, assume that $v\in R^n$. Then there is an element $A\in{\rm SL}_n(R)$ such that $A\cdot v\in Re_1\oplus Re_2.$
\end{lemma}

\begin{lemma}\label{stability_sln_prep_step1}
Let $R$ be a Dedekind domain and let $I$ be a non-zero ideal in $R.$ Further, let $n\geq 2$ and $v\in R^n$ be given such that $v\equiv e_1\text{ mod }I$ and such that the entries of $v$ generate $R$ as an ideal. Then there is a product $X$ of six elements of $Z(I,A_n)$ such that $X\cdot v=e_1.$ 
\end{lemma}

\begin{lemma}\label{fixing_first_column_symplectic_step1}
Let $R$ be a Dedekind domain, $n\geq 2$ and let $A\in{\rm Sp}_{2n}(R)$ be given. Then the conjugacy class of $A$ in ${\rm Sp}_{2n}(R)$ contains an element $A'=(a_{ij}')_{1\leq i,j\leq 2n}$ such that the first $n$ entries $a'_{11},\dots,a'_{n1}$ of the first column of $A'$ generate $R$ as an ideal.  
\end{lemma}

\begin{lemma}\label{fixing_first_column_symplectic_step2}
Let $R$ be a Dedekind domain and let $I$ be a non-zero ideal in $R.$ Further, let $n\geq 2$ be given and $A\in C(C_n,I)$ be given. Then up to multiplication with $15$ elements of $Z(I,C_n)$, the first and $n+1.$th columns of $A$ and the first and $n+1$.th row of $A$ can be assumed to be $e_1, e_{n+1}, e_1^T$ and $e_{n+1}^T$ respectively. 
\end{lemma}

\begin{proof}
We first note that by applying Lemma~\ref{fixing_first_column_symplectic_step1}, we can assume that the first $n$ entries $a_{11},\dots,a_{n1}$ of the first column of $A$ already generate the ring $R.$ The vector $v:=(a_{11},\dots,a_{n1})^T$ further satisfies $v\equiv e_1^T\text{ 
 mod } I$. Thus applying Lemma~\ref{stability_sln_prep_step1} to the ${\rm SL}_n(R)$ contained in ${\rm Sp}_{2n}(R)$ via
\[
{\rm SL}_n(R)\to{\rm Sp}_{2n}(R),M\mapsto
\begin{pmatrix}
M & 0_n\\
0_n & M^{-T}.
\end{pmatrix}
\]
and $v$ we can find an element $X$ which is a product of six elements of $Z(I,C_n)$ such that $B:=X\cdot A$ has first column with its first $n$ entries equal to $e_1^T\in R^n.$ Then consider the product
\begin{align*}
C:=&(I_{2n}-b_{2n,1}(e_{n+1,n}+e_{2n,1}))\cdots(I_{2n}-b_{n+3,1}(e_{n+1,3}+e_{n+3,1}))\cdot(I_{2n}-b_{n+2,1}(e_{n+1,2}+e_{n+2,1}))\\
&\cdot(I_{2n}-b_{n+1,1}e_{n+1,1})\cdot B
\end{align*}
and note that its first column is precisely $e_1\in R^{2n}.$
Next, we will show that the product 
\[
(I_{2n}-b_{2n,1}(e_{n+1,n}+e_{2n,1}))\cdots(I_{2n}-b_{n+3,1}(e_{n+1,3}+e_{n+3,1}))\cdot(I_{2n}-b_{n+2,1}(e_{n+1,2}+e_{n+2,1}))
\]
is a product of two elements from $Z(I,C_n)$. To this end, consider the subgroup ${\rm SL}_{n-1}(R)$ of ${\rm Sp}_{2n}(R)$ given by the map
\[
j:{\rm SL}_{n-1}(R)\to{\rm Sp}_{2n}(R),M\mapsto
\begin{pmatrix}
1 & & &\\
& M^{-T} & & \\
& & 1 & \\
& & & M.
\end{pmatrix}.
\]
Note that for an element $v\in R^{n-1}$ and the matrix $Y(v)\in R^{2n\times 2n}$ of the form 
\[
Y(v)=I_{2n}+\sum_{i=1}^{n-1} v_i\cdot (e_{n+i+1,1}+e_{n+1,i+1})
\]
conjugation of $Y(v)$ by $j(M)$ acts on $Y(v)$ by the standard operation of ${\rm SL}_{n-1}(R)$ on $R^{n-1},$ that is $j(M)Y(v)j(M)^{-1}=Y(Mv)$ holds for all $M\in{\rm SL}_{n-1}(R)$ and $v\in R^{n-1}.$ Hence it suffices to show that for any vector $v\in R^{n-1}$ there is some $M\in{\rm SL}_{n-1}(R)$ with $M\cdot v\in Re_1\oplus Re_2.$ But this is precisely the claim of Lemma~\ref{stability_relative_ideal}. So after multiplying $A$ with $9$ elements of $Z(I,C_n),$ we can assume that the first column of $C$ is $e_1.$ 

The $1$ in the first column of $C$ can now be used to clear out the first row of $C$: Picking suitable elements $x_2,\dots,x_n,y_{n+1},\dots,y_{2n}\in I$ one obtains that the product
\begin{align*}
C&\cdot(I_{2n}+x_2(e_{12}-e_{n+2,n+1}))\cdot(I_{2n}+x_3(e_{13}-e_{n+3,n+1}))\cdots(I_{2n}+x_n(e_{1,n}-e_{2n,n+1}))\cdot(I_{2n}+y_{n+1}e_{1,n+1})\\
&\cdot (I_{2n}+y_{n+2}(e_{1,n+2}+e_{2,n+1}))\cdot (I_{2n}+y_{n+3}(e_{1,n+3}+e_{3,n+1}))\cdots(I_{2n}+y_{2n}(e_{1,2n}+e_{n,n+1}))
\end{align*}
called $D$, has the first row equal to $e_1^T.$ However, as before, one obtains using Lemma~\ref{stability_relative_ideal} that the two products 
\begin{align*}
&(I_{2n}+x_2(e_{12}-e_{n+2,n+1}))\cdot(I_{2n}+x_3(e_{13}-e_{n+3,n+1}))\cdots(I_{2n}+x_n(e_{1,n}-e_{2n,n+1}))
\text{ and }\\
&(I_{2n}+y_{n+2}(e_{1,n+2}+e_{2,n+1}))\cdot (I_{2n}+y_{n+3}(e_{1,n+3}+e_{3,n+1}))\cdots(I_{2n}+y_{2n}(e_{1,2n}+e_{n,n+1}))
\end{align*}
are products of two elements of $Z(I,C_n)$ respectively. So in total after multiplying $A$ with $9+4+1=15$ elements of $Z(I,C_n),$ we have reduced $A$ to an element $D$ with first column $e_1$ and first row $e_1^T.$ But this actually puts $D$ into the required form: The columns of $D$ form a symplectic basis wrt the symplectic form given by the matrix $J\in R^{(2n)\times (2n)}.$ Hence note that one has for the k.th column $v_k$ of $D$ that $0=e_1^T\cdot J\cdot v_k=d_{n+1,k}$ if $k\neq n+1.$ Hence the $n+1.$th row of $D$ is $d_{n+1,n+1}e_{n+1}^T=e_{n+1}^T.$ 

However, $D$ is a symplectic matrix and so for the $n+1.$st column
\[
\begin{pmatrix}
w_1\\
w_2
\end{pmatrix}.
\]
of $D$ with $w_1,w_2\in R^n$, the first row of its inverse $E$ is given by $(w_2^T,-w_1^T)$. But using that $D$ is symplectic the $n+1$.th column of $E$ is given by $e_{n+1}.$ Then as before, one can use the fact that the columns of $E$ form a symplectic basis wrt $J$ to deduce that the first row of $E$ is equal to $e_1$. Hence the $n+1.$th column of $D$ is $e_{n+1}$ and we are done. 
\end{proof}

\end{document}
